% 图契约：本题的源、边界和推导逐项对应；箭头符号按图注解释。
\documentclass[tikz,border=7pt]{standalone}
\usepackage{fontspec,xeCJK,amsmath,amssymb}
\IfFontExistsTF{Times New Roman}{\setmainfont{Times New Roman}}{\setmainfont{TeX Gyre Termes}}
\IfFontExistsTF{Heiti SC}{\setCJKmainfont[AutoFakeBold=1.5]{Heiti SC}}{\setCJKmainfont[AutoFakeBold=1.5]{FandolHei-Regular}}
\usetikzlibrary{arrows.meta,bending,calc,patterns,positioning}
\definecolor{ink}{HTML}{243745}
\definecolor{blue}{HTML}{275DAB}
\definecolor{green}{HTML}{177C63}
\definecolor{red}{HTML}{BD4935}
\definecolor{orange}{HTML}{AA6C16}
\definecolor{purple}{HTML}{715393}
\tikzset{every node/.style={font=\fontsize{9.3}{12.5}\selectfont,text=ink},
axis/.style={-{Stealth[length=4pt]},line width=.65pt,draw=ink!65},
vec/.style={-{Stealth[length=5pt]},line width=1pt,draw=blue},
short/.style={-{Stealth[length=3pt]},line width=.8pt},
guide/.style={dashed,draw=ink!40,line width=.5pt},
card/.style={rounded corners=2pt,draw=blue!65,fill=blue!4,align=center,text width=6.5cm,inner xsep=9pt,inner ysep=7pt},
title/.style={font=\fontsize{10.5}{14}\selectfont\bfseries,text width=15cm,align=center},
label/.style={font=\fontsize{8.4}{11}\selectfont,align=center},
note/.style={label,text width=6.2cm}}
\begin{document}
\begin{tikzpicture}[x=1cm,y=1cm]
\path[use as bounding box] (-8,-3.5) rectangle(8,3.5);
\node[title] at (0,3.15) {M3-009\quad 没有自由电荷时的非零电位移};
\begin{scope}[shift={(-4,0)}]
\draw[green,dashed,thick] (-1.1,-1.15) rectangle(1.1,1.15);
\node[green] at(0,1.5){闭合面 $S$};
\foreach \y in {-.7,0,.7}{\draw[vec] (-2.1,\y)--(2.1,\y);}
\draw[vec,orange] (-1.1,-1.15)--(-1.7,-1.15) node[below] {$\mathbf n_L$};
\draw[vec,orange] (1.1,-1.15)--(1.7,-1.15) node[below] {$\mathbf n_R$};
\draw[ink,thick] (-2.3,-1.5)--(-2.3,1.5) (2.3,-1.5)--(2.3,1.5);
\node at(-2.3,1.9){$+\sigma_f$};\node at(2.3,1.9){$-\sigma_f$};
\node[fill=white] at(0,-.35){$q_{f,S}=0$};\node[label] at(0,-2.1){两极板均位于 $S$ 外；箭头为 $\mathbf D$};\end{scope}
\node[card,draw=blue!65,fill=blue!4] at (3.9,1.1) {$q_{f,S}=0\ \Longrightarrow\ \Phi_D=0$};
\node[card,draw=green!65,fill=green!4] at (3.9,-1.1) {外部自由电荷仍可产生\\非零且穿过 $S$ 的 $\mathbf D$};
\end{tikzpicture}\end{document}